(1) 取正径向，圆柱高斯面给 $D_r2\pi r=\lambda$，故 $D_r=\lambda/(2\pi r)$。在线性各向同性介质内 $E_r=D_r/\varepsilon(r)=\lambda/[2\pi\varepsilon_0r(1+\alpha r^2/a^2)]$，这里场强与位移矢量不可在非均匀介质中混同。
(2) 设 $F=\ln(b/a)-\frac12\ln[(1+\alpha b^2/a^2)/(1+\alpha)]$。由 $V=\int_a^bE_r\,dr=\lambda F/(2\pi\varepsilon_0)$，得到 $C'=\lambda/V=2\pi\varepsilon_0/F$。由于 $\varepsilon(r)>0$，实际积分 $F=\int_a^b dr/[r(1+\alpha r^2/a^2)]$ 严格正。
(3) 每长度场能 $U'=\frac12\int_a^b\mathbf E\cdot\mathbf D\,2\pi r\,dr=\lambda^2F/(4\pi\varepsilon_0)=\lambda^2/(2C')$，与充电功积分一致。增大 $\alpha$ 将正区域中的 $\varepsilon(r)$ 增大，积分 $F$ 减小，因此 $C'$ 单调增大，而非由电势差符号猜测单调性。
(4) 用 $\ln(1+z)=z+O(z^2)$，有 $F\to\ln(b/a)$、$C'\to2\pi\varepsilon_0/\ln(b/a)$。两个对数的自变量都是长度比或介电常数比，不可对有量纲的 $r$ 单独取对数。要求全部 $r\in[a,b]$ 上 $1+\alpha r^2/a^2>0$；题设末端正性结合 $\alpha$ 的单调结构即可保证。
